% !TeX TS-program = xelatex
\documentclass{resume}
\ResumeName{陈沫樵}
\usepackage{graphicx}
\usepackage{tikz}
\usepackage{enumitem} % 用于优化列表间距

\begin{document}

\ResumeContacts{
  (+86) 186\-1195\-1800,%
  \ResumeUrl{mailto:mochiaochen@gmail.com}{mochiaochen@gmail.com},
  \ResumeUrl{github.com/mochiao}{github.com/MochiaoChen}
}

\ResumeTitle

\section{教育经历}
\ResumeItem
[对外经济贸易大学 | 本科]
{对外经济贸易大学}
[金融学（鸿儒金融人才培养实验班），中国金融学院 | 经济学学士]
[2024.09—2028.06]

\section{技术能力}
\begin{itemize}[parsep=2pt]
  \item \textbf{编程与开发}: 熟练掌握 \textbf{Python}，擅长数据分析；熟悉\textbf{C++}、\textbf{SQL}, \textbf{HTML/JavaScript} 及 \LaTeX{}；了解 \textbf{Vibe Coding}，熟练掌握\textbf{大语言模型应用},熟练掌握\textbf{Prompt Engineering}。
  \item \textbf{数理与分析}: 掌握\textbf{机器学习}基础与\textbf{数理统计建模}；能够独立完成\textbf{数据清洗}、\textbf{可视化分析}及\textbf{行业研究简报}撰写。
  \item \textbf{语言能力}: 英语 \textbf{C1} 水平 (CET-4 \textbf{626}分)，具备优秀的英文专业文献阅读与学术写作能力。英语可作为\textbf{工作语言}
\end{itemize}

\section{实习经历}

\ResumeItem{\textbf{投资银行部实习生}}
[中信证券]
[2026.01-至今]

\begin{itemize}

\item \textbf{行业洞察}: 深度参与某具身智能企业港股 IPO 项目，负责招股书中“业务”章节的撰写。通过对 AI 算法、感知系统及机器人硬件链条的梳理，完成核心业务逻辑分析与多家同业竞对研究。
\item \textbf{流程支持}: 负责协助 M111 表格填报及港交所回函撰写，确保上市申请流程符合港交所合规标准。

\end{itemize}

\ResumeItem{\textbf{投资银行部实习生}}  
[长城证券]  
[2025.09—2026.03]  
\begin{itemize}  
  \item \textbf{项目承做}: 深度参与\textbf{公司债}项目的 \textbf{尽职调查}，协助撰写立项报告与申报材料，整理底稿并完成归档。    
  \item \textbf{数据分析}: 运用 \textbf{Wind}, \textbf{iFind} ，\textbf{Excel}及 \textbf{Python} 进行市场数据清洗与同行业分析；
  \item \textbf{AI应用}：构建一系列用脚本进行文档处理、并使用CLI agent 自动化处理部分工作，有效提升工作效率。  
\end{itemize}

\ResumeItem{\textbf{AI科技内容运营实习生}}  
[云杉资本]  
[2025.03—2025.09]  
\begin{itemize}  
  \item \textbf{行业研究}: 聚焦 \textbf{AI} 与 \textbf{消费} 赛道，独立撰写多篇行业分析简报，拆解国内外 AI 初创项目的商业模式与技术路径。  
  \item \textbf{内容策划}: 建立 AI 赛道投资视角话题库，策划并输出深度分析文章，辅助投资团队进行市场趋势研判。  
\end{itemize}

\section{项目经历}

\ResumeItem{\textbf{“工行杯”全国大学生金融科技创新大赛 | 全国三等奖}}
[核心开发者 / Tech Lead]
[2025.11]
\begin{itemize}
\item \textbf{项目概况}: 开发“鉴真盾”:基于多模态大模型的金融反欺诈人像风险识别系统，解决传统信审中黑产攻击等痛点。
\item \textbf{技术实现}: 设计 \textbf{Agent} 协作架构，通过 \textbf{Prompt Engineering} 精准识别六大类欺诈特征；负责全栈开发。
\item \textbf{量化成果}: 建立了中介欺诈特征库，模型在测试集准确率达 \textbf{90\%}，实现风险量化评分机制。
\end{itemize}



\ResumeItem{\textbf{“挑战杯”首都大学生课外学术科技作品竞赛 | 北京市三等奖}}  
[项目负责人 / PM]  
[2025.05]  
\begin{itemize}  
  \item \textbf{产品设计}: 负责“心桥”Agent（AI 心理健康伙伴）的全流程落地，包括\textbf{产品PRD设计}与\textbf{商业计划书}撰写。  
  \item \textbf{功能开发}: 结合 \textbf{RAG} 技术与心理学咨询知识库，设计具备情绪识别与心理辅导、陪伴功能的 AI 助手。
\end{itemize}

\ResumeItem{\textbf{全国大学生数学建模竞赛 | 北京赛区一等奖}}
[核心成员]
[2025.10]
\begin{itemize}
  \item \textbf{项目概况}: 针对无创产前检测 (NIPT) 时点优化与异常判定问题，基于临床数据构建多维度机器学习模型体系。
  \item \textbf{数据分析与建模}: 负责完成原始数据缺失值插补、异常值剔除等预处理工作；协助建模手编程。
  \item \textbf{成果产出}: 负责撰写项目核心摘要，并使用 Python 绘制数据分布、特征重要性、ROC 曲线等多维可视化图表。
\end{itemize}

\ResumeItem{\textbf{党建智能 Agent 系统 | 学院合作项目}}  
[Prompt 开发者]  
[2025.03]  
\begin{itemize}  
  \item \textbf{架构搭建}: 基于 \textbf{MetaGPT} 框架构建多智能体系统，集成 \textbf{UltraRAG} 向量检索模块。  
  \item \textbf{效果优化}: 设计多角色 Prompt 模板，实现了高精度的政策解读、报告生成、党建学习等功能。  
\end{itemize}

\section{所获奖项}
\begin{itemize}[leftmargin=1.5em, parsep=0pt, itemsep=2pt]
  \item 2025年北京市大学生人文知识竞赛 \textbf{个人赛一等奖} \hfill 2025.11
  \item 2025年“工行杯”全国大学生金融科技创新大赛 \textbf{全国三等奖} \hfill 2025.10
  \item 2025年全国大学生数学建模竞赛 \textbf{北京赛区一等奖} \hfill 2025.10
  \item “青创北京”2025年“挑战杯”首都大学生课外学术科技作品竞赛 \textbf{三等奖} (负责人) \hfill 2025.05
  \item 2025年(第十一届)全国大学生统计建模大赛 \textbf{三等奖} \hfill 2025.05
  \item 第八届对外经济贸易大学信息检索大赛 \textbf{二等奖} \hfill 2024.11
  \item 第37届中国化学奥林匹克 \textbf{陕西省一等奖} \hfill 2023.11
\end{itemize}

\end{document}